\section{Entropy production for specified matrix ensembles}\label{sec:flow}

We use the classical-to-free Fisher comparison of
Jekel and Pi \cite[Proposition~6.1]{JP},
keeping the terminal ensemble in the conclusion. The purpose is to
transfer a later entropy bound for one evolved matrix law back to the
initial microstates entropy. We prove the comparison for the exact
laws that will enter the structural event in the next section.

\begin{lemma}\label{lem:flow}
For the fixed tuple $X=V_h$, fix a cutoff $R>0$ with
$\chi_R(X)>-\infty$. There are dimensions tending to infinity and
random tuples $A^{(d)}$, uniform on increasingly accurate
positive-volume microstate sets at cutoff $R$, such that, for independent
standard GUE tuples,
\begin{equation}\label{eq:flowlower}
 \limsup_d H_d(A^{(d)}+\sqrt t G^{(d)})
 \geq\chi_R(X)+\frac12\int_0^t\Phi^*(V_{h+s})\,ds.
\end{equation}
The left side is the limsup of the normalized entropies of these particular ensembles.
\end{lemma}

\begin{proof}
Every uniform law on a positive-volume microstate set at cutoff $R$
has $\sum_i\E\tr_d(A_i^2)\leq nR^2$. Comparison with independent
GUE matrices of variance $R^2$ bounds its normalized entropy above by
$\frac n2\log(2\pi eR^2)$. Thus $\chi_R(X)$ is a finite real number.
Choose a new increasing sequence of dimensions $d_k$, $k\geq1$, such that the normalized log-volume
at parameters $(m,\varepsilon)=(k,1/k)$ is at least
$\chi_R(X)-1/k$. Such arbitrarily large dimensions exist by the defining
limsup, and the selected sets have positive volume. The pairs $(k,1/k)$
are cofinal among all moment accuracies. Their uniform laws satisfy
\begin{equation}\label{eq:initialentropy}
 \liminf_d H_d(A^{(d)})\geq\chi_R(X).
\end{equation}
Here and below $d$ ranges over the selected dimensions $d_k$.
Every tuple in the support of $A^{(d)}$ has the same fixed cutoff and
increasingly accurate moments of $X$. The selection does not depend on
the subsequent noise time or polynomial tests.

Put $Z^{(d)}(s)=A^{(d)}+\sqrt s G^{(d)}$, and let $f_s$ be its density
for $s>0$. In the $nd^2$ real Hilbert--Schmidt coordinates, the added
Gaussian has covariance $(s/d)I$. Thus
$\partial_s f_s=(1/(2d))\Delta f_s$. Differentiating the entropy and
integrating by parts gives the normalized de Bruijn identity
\begin{equation}\label{eq:matrix-debruijn}
 \frac{d}{ds}H_d(Z^{(d)}(s))=\tfrac12 I_d(s),\qquad
 I_d(s)=\frac1{d^3}\int\|\nabla\log f_s\|^2 f_s
       =\sum_i\E\tr_d(F_i^2),
\end{equation}
where $F_i=-(1/d)\nabla_i\log f_s$ evaluated at $Z^{(d)}(s)$.
Gaussian convolution of a compactly supported distribution justifies
differentiation and integration by parts on compact positive-time
intervals. Indeed, derivatives have Gaussian decay, a Gaussian lower
bound on $f_s$ bounds $|\log f_s(x)|$ by $C(1+\|x\|^2)$, and
\[
 \nabla\log f_s(x)
 =-\frac d s\big(x-\E[A^{(d)}\mid Z^{(d)}(s)=x]\big)
\]
has at most linear growth. Here the norm and gradient are Euclidean,
and the conditional expectation is bounded at each fixed dimension.
In particular the entropy and Fisher information are finite at positive
times.
Entropy concavity under mixtures of translations gives
$H_d(Z^{(d)}(\varepsilon))\geq H_d(A^{(d)})$. Integrating first from
$\varepsilon$ to $t$ and then letting $\varepsilon\downarrow0$ proves
\begin{equation}\label{eq:classgain}
 H_d(Z^{(d)}(t))\geq H_d(A^{(d)})+\tfrac12\int_0^t I_d(s)\,ds.
\end{equation}
No continuity of entropy at zero is needed.

For fixed $s>0$, integration by parts against any complex
noncommutative polynomial $P$ gives
\begin{equation}\label{eq:matrixibp}
 \E\tr_d(F_i P(Z^{(d)}(s)))
 =\E(\tr_d\otimes\tr_d)((\partial_iP)(Z^{(d)}(s))).
\end{equation}
To verify this identity including its normalization, take an orthonormal
Hermitian basis $(E_\alpha)$ for unnormalized Hilbert--Schmidt inner
product. It satisfies
$\sum_\alpha E_\alpha D E_\alpha=\Tr(D)1$ for every complex matrix $D$.
Scalar integration by parts, separately on real and imaginary parts,
turns the left side of \eqref{eq:matrixibp} into
\[
 d^{-2}\sum_\alpha\E\Tr(E_\alpha D_iP(Z^{(d)}(s))[E_\alpha]).
\]
Here $D_iP(Z)[E_\alpha]$ is the directional derivative in the $i$th
matrix coordinate along $E_\alpha$.
At an occurrence $P=LZ_iR$ in a monomial, its contribution is
\[
 d^{-2}\sum_\alpha\Tr(E_\alpha L E_\alpha R)
 =d^{-2}\Tr(L)\Tr(R)=\tr_d(L)\tr_d(R).
\]
Summing over occurrences proves \eqref{eq:matrixibp} with the prefix
and suffix in the same order as the free difference quotient.

For self-adjoint polynomial tests $P_i$, completion of squares gives
\[
 I_d(s)\geq\sum_i\E\tr_d\big(2F_iP_i(Z^{(d)}(s))
                                  -P_i(Z^{(d)}(s))^2\big).
\]
First use \eqref{eq:matrixibp} to replace every term involving the score
by a product of normalized polynomial traces. Lemma~\ref{lem:matrix-convergence}
then gives the limit of the right side as
\[
 \sum_i\left[2(\tau\otimes\tau)((\partial_iP_i)(V_{h+s}))
                         -\tau(P_i(V_{h+s})^2)\right].
\]
The limiting law is $V_{h+s}$ because the new semicircular noise is
free from $X=V_h$. Conjugate variables exist by Lemma~\ref{lem:wick}.
Taking the supremum over these tests in Lemma~\ref{lem:variational} gives
\begin{equation}\label{eq:matrix-free-fisher}
 \liminf_d I_d(s)\geq\Phi^*(V_{h+s}).
\end{equation}
This holds for every fixed $s>0$ along the same selected dimensions.
The moments of $V_{h+s}=C+\sqrt{h+s}S$ are continuous in $s$, so
Lemma~\ref{lem:variational} also gives measurability of the free Fisher
information. Applying Fatou's lemma to the nonnegative measurable
functions $I_d$, and then using \eqref{eq:initialentropy} and
\eqref{eq:classgain}, yields
\[
 \liminf_d H_d(Z^{(d)}(t))
 \geq\chi_R(X)+\frac12\int_0^t\Phi^*(V_{h+s})\,ds.
\]
This proves \eqref{eq:flowlower}, even with liminf in place of limsup.
The free Fisher integral is finite since
$\Phi^*(V_{h+s})\leq n/(h+s)$ and $h>0$. Throughout the proof the
terminal law is that of the same $A^{(d)}+\sqrt tG^{(d)}$ selected above;
no entropy-maximizing law at the terminal distribution has been chosen.
\end{proof}
